\section*{Part VIII: International and Multientity Governance}
\addcontentsline{toc}{section}{Part VIII: International and Multientity Governance}

\subsection{53. Cross-border entities are not just bank accounts}

A foreign branch, subsidiary, agent, reseller, academic partner and foundation create different permanent-establishment, tax, employment, transfer-pricing, data and governance risks. Before establishing foreign presence, map who signs contracts, where services are performed, who owns IP, who employs people, where revenue is recognized and where decision makers act. A global website does not by itself require foreign incorporation, but ongoing local operations may.

Internationalization should be staged and supported by local legal/tax advisers in the actual jurisdictions.

\subsection{54. Intercompany agreements and transfer pricing}

Where related entities supply software, research, curriculum, property, management services or licenses, maintain lawful intercompany agreements, allocation methodologies and tax/transfer-pricing records where applicable. A for-profit company should not siphon value from a donor-supported research entity through inflated management fees. Nor should a foreign subsidiary receive unpriced core IP without reviewing tax and shareholder consequences.

Keep the \textbf{economic purpose} of each transfer inspectable. Internal accounting convenience is not evidence of lawful pricing.

\subsection{55. International board and delegation design}

Specify which local manager can price, hire, sign customer contracts, select partners, approve data handling and incur capital commitments; which matters require central approval; and how local compliance issues reach the parent board. Overcentralization slows customers; undercontrolled country teams can create contractual and regulatory liabilities that appear months later.

A globally coherent institution needs consistent integrity standards and locally competent operating authority.
